\section{Deployment and Apache Jena Fuseki}
\label{sec:fuseki}
This section documents how the validated RDF release is published and queried. It is useful
to separate the URI design from the runtime: a URI identifies a resource, while Fuseki stores
and queries triples about that resource. The same graph can therefore be published as static
RDF and loaded into Fuseki without changing its identifiers.
Fuseki is a SPARQL server, not a URI minting service~\cite{fuseki}.
The transformer assigns canonical identifiers before data enters any store.
The same university IRI remains in RDF whether queried by RDFLib, Oxigraph or Fuseki.

\Needspace{15\baselineskip}
\subsection{Identity, representation and storage}
\begin{code}
# RDF subject identifying a university:
https://pham-ng.github.io/Vietnam-University-Knowledge-Graph-ver2/resource/university/dai-hoc-bach-khoa-ha-noi
# Local query endpoint:
http://localhost:3030/vnedu/sparql
# Graph Store Protocol target:
http://localhost:3030/vnedu/data?default
\end{code}
Uploading RDF does not make its canonical IRIs dereferenceable automatically.
The publisher independently creates HTML, Turtle and JSON-LD descriptions. Ontology terms
retain hash IRIs; entity and observation identifiers use slash paths.
Immutable 2.1 and 2.2 version artifacts are generated under \term{ontology/}.
Dynamic version routes support RDF content negotiation and return 404 for unknown versions.

GitHub Pages remains a separate HTTP origin from Flask. Its static representation links
and embedded JSON-LD do not implement Flask's negotiation, nor the entity-to-document
303 pattern discussed in \emph{Cool URIs}~\cite{uris}. Adding Fuseki cannot fix those origin-level
behaviors. Version artifacts are implemented and tested in the release; availability from
each public origin remains an HTTP-deployment concern rather than an ontology property.

\begin{figure}[htbp]\centering
\resizebox{\textwidth}{!}{\begin{tikzpicture}
\node[auditbox,text width=4cm] (gold) at (0,0) {Validated union\\vnedu-all.ttl};
\node[auditbox,text width=4.3cm] (static) at (8.5,0) {Static publisher / GitHub Pages\\Canonical entity HTML + RDF documents};
\node[auditbox,text width=4cm] (rdf) at (0,-2.5) {RDFLib child processes\\Default Render execution path};
\node[auditbox,text width=4.3cm] (flask) at (8.5,-2.5) {Flask /sparql and /resource/...\\Query guard + response handling};
\node[auditbox,text width=4cm] (fuseki) at (0,-5.3) {Optional Fuseki /vnedu\\In-memory file OR persistent TDB2};
\node[auditbox,text width=4.3cm] (client) at (8.5,-5.3) {Client: browser or RDF consumer\\Chooses a document or query endpoint};
\draw[auditarr] (gold)--node[auditlabel,above]{build}(static);
\draw[auditarr] (gold)--node[auditlabel,right]{read}(rdf);
\draw[auditarr] (flask)--node[auditlabel,above]{local query}(rdf);
\draw[auditarr,dashed] (gold.west) -- ++(-0.7,0) |- node[auditlabel,pos=0.65,left]{file load /\\explicit PUT}(fuseki.west);
\draw[auditarr,dashed] (flask.south west)--node[auditlabel,sloped,above]{optional proxy}(fuseki.north east);
\draw[auditarr] (client)--node[auditlabel,right]{HTTP query}(flask);
\draw[auditarr] (client.east)--++(0.7,0)|-node[auditlabel,pos=0.45,right]{HTTP GET}(static.east);
\end{tikzpicture}
}
\caption[URI publication and optional Fuseki execution]{Publication and storage paths. The optional Fuseki
branch has verified Linux and Windows acceptance runs; Render still explicitly selects local RDFLib.
Storage does not replace the canonical identifiers.}\label{fig:deployment}
\end{figure}

\subsection{Safe loading and reproducible runtime}
The revised loader requires a successfully validated release before network writes.
It targets an existing explicitly named dataset rather than creating one implicitly.
An empty graph can be loaded directly. Replacement of a nonempty default graph requires
both \term{--replace} and a new \term{--backup} path; the backup is created exclusively,
so an older backup cannot be silently overwritten. HTTP redirects and error responses are
rejected. After PUT, the loader retrieves the entire default graph and checks RDF graph
isomorphism, including blank nodes, rather than relying on a triple count.

Graph Store PUT replaces the selected graph~\cite{gsp}; named graphs are untouched by
this loader. Credentials are read from environment variables, not command-line passwords.
A failed comparison preserves the backup for operator recovery; it is not a distributed
transaction or automatic rollback. Operators must stop concurrent writers during a release:
a backup followed by PUT does not provide compare-and-swap concurrency control.

The runtime setup pins Apache Jena Fuseki 6.2.0 and verifies its published SHA-512 digest.
The local launcher requires Java 21, validates the release, binds to loopback, and enables
only query access with a configured timeout. Its file mode reloads the graph at startup;
the separate acceptance harness explicitly tests persistent TDB2.
The optional legacy Docker recipe now requires an operator-selected image and password,
and binds port 3030 only on loopback. That Docker configuration was not the tested runtime.

\subsection{Production-like security and performance acceptance}
The additional \term{audit/production\_readiness.py} test starts the same Waitress WSGI
entry point used by \term{app/server.py --prod} with the read-only local backend. Under
normal bounded concurrency, all eight constant-time health requests returned HTTP 200; the
final 9 October 2026 rerun measured 12.45 ms at p50, 15.89 ms at p95 and 15.89 ms maximum
on the Windows review machine. An intentional simultaneous eight-request SPARQL overload
against two query workers returned two HTTP 200 responses and six HTTP 503 responses,
demonstrating bounded back-pressure rather than unbounded query growth. A parser lock also
prevents concurrent cold requests from racing RDFLib's lazy SPARQL grammar initialization.
Security probes rejected a path-injection
attempt (404), an unapproved SSRF \term{SERVICE} target (400), and an oversized query (413).
The response headers now include \term{nosniff}, \term{SAMEORIGIN}, strict referrer-policy,
\term{Permissions-Policy}, CSP on HTML, and an \term{X-Request-ID}; HSTS is enabled when
TLS is detected or explicitly forced behind a trusted TLS terminator.

These numbers are a reproducible loopback acceptance sample, not a capacity SLO, penetration
test, TLS/reverse-proxy audit, cloud IAM review, or production security certification. The
bounded persistent RDFLib worker deliberately trades throughput for isolation; a real
deployment still needs external rate limiting, resource quotas, TLS termination, monitoring,
authentication policy, and a load test against its selected Fuseki/WSGI topology. The full
machine-readable evidence is \term{data/reports/production-readiness.json}.

\subsection{Production hardening profile and release gates}
The review branch adds a deployable reference profile in \term{deploy/} and records its
scope in \term{docs/production-readiness.md}. The intended production path is:
\begin{code}
Internet -> Caddy/Nginx TLS -> Waitress application -> private Fuseki -> encrypted TDB2 volume
\end{code}
The application layer is read-only for SPARQL, limits request bodies and serialized result
size, applies an optional per-client rate limit, restricts CORS to an explicit allow-list,
and emits request identifiers for incident correlation. The reference Compose file also
uses a non-root application user, a read-only container filesystem, a private network,
dropped Linux capabilities and a no-new-privileges flag. Fuseki image digests and admin
credentials are deliberately operator-supplied; the public reverse proxy must never expose
Fuseki administration or update endpoints.

Some application-layer controls are active on the public Render service, but this is not
evidence that the complete reverse-proxy/Fuseki/TDB2 profile has been deployed. Before a
production claim for that topology, the deployed Fuseki path
must be tested for sustained and burst requests, p50/p95/p99 latency, CPU/memory/disk
growth, query cancellation, rolling restart and backup restoration. The current local
RDFLib p95 is therefore reported as an isolation/back-pressure observation, not as a
capacity target.

\begin{table}[htbp]\centering\small
\caption{Production release gates that remain separate from OWL and SHACL conformance.}
\label{tab:production-gates}
\begin{tabularx}{\textwidth}{p{4.2cm}X}\toprule
Gate & Required evidence before an authoritative production claim \\\midrule
Security deployment & Public TLS termination, trusted proxy configuration, firewall/IAM policy, secret rotation and external rate limits. \\
High-load performance & Deployed Fuseki benchmark with representative SELECT/CONSTRUCT queries, p50/p95/p99, saturation behavior and resource measurements. \\
Recovery & The immutable publication ZIP was restored and hash/parse checked locally; runtime database snapshots, off-host retention and rolling restore remain operator gates. \\
Factual accuracy & Independent adjudication of the five founding-year conflicts and seven unresolved values, plus a held-out identity-link evaluation. \\
Temporal validity & Claim-level provenance and validity intervals for leadership, enrolment and legal status; undated snapshots must not be presented as current. \\
Source authority & Official-register or institution-owner evidence for records currently supported only by encyclopedic or inferred sources. \\
\bottomrule\end{tabularx}
\end{table}

The release can pass the OWL 2 RL profile gate, SHACL validation and regression tests while
still failing these operational and factual gates. Keeping those gates separate is essential:
OWL entailment is not a truth certificate, and a five-star Linked Data publication is not a
five-star security or accuracy certification.

\subsection{Executed end-to-end acceptance evidence}
The primary reproducible test ran on Linux GitHub Actions with Fuseki 6.2.0 and Temurin Java 21.0.12.1,
on 7 October 2026 at 01:50 UTC. The latest Windows run began on 9 October 2026 at
03:50:45 UTC and used the same pinned Fuseki JAR, Java 21.0.8 and the selector/socket
workaround against the current 87,183-triple release. An accidental invocation through the
Java 8 executable failed before startup because Fuseki 6.2 requires class-file version 65;
rerunning with the explicitly selected JDK 21 completed successfully. This is useful runtime
selection evidence rather than a semantic or data failure.
The input hash, JAR hash, runtime details and assertions are archived in
\term{data/reports/fuseki-e2e-linux.json} and the Windows run in
\term{data/reports/fuseki-e2e-windows.json}; the corresponding Linux
\href{https://github.com/pham-ng/Vietnam-University-Knowledge-Graph-ver2/actions/runs/37559070543}{CI run}
completed successfully.

\begin{table}[htbp]\centering\small
\caption{Measured Fuseki/TDB2 acceptance results.}\label{tab:fuseki-status}
\begin{tabularx}{\textwidth}{P{6.5cm}X}\toprule
Check & Observed result \\\midrule
Load complete validated serving graph & 87,183 triples; graph-isomorphic to local RDF. \\
Attempt unapproved replacement & Refused before PUT. \\
Approved replacement with backup & Backup graph-isomorphic to previous contents. \\
Query parity with RDFLib & Count: 1 row; class totals: 68 rows; programs: 78 rows; explicitly dated leadership: 0 rows. \\
Stop and restart with the same TDB2 directory & Entire default graph remains isomorphic. \\
Restart without update permission & SPARQL Update and Graph Store PUT both return HTTP 405. \\
Read after rejected writes & Graph unchanged and isomorphic. \\
\bottomrule\end{tabularx}
\end{table}

The test uses a newly created disposable directory and ephemeral loopback port; no existing
user dataset is overwritten. Zero dated-leadership rows is an honest reflection of missing
date evidence, not a failed lookup to be filled with invented dates.

On Windows, the first attempt exposed a JDK 21 selector/socket initialization failure.
The repaired launcher now sets the selector provider and short Unix-domain socket directory
through inherited Java options (avoiding PowerShell 5.1 argument concatenation); the same
disposable TDB2 acceptance test then passed all checks, and a direct query-only launcher
smoke test returned HTTP 200. This proves local Windows operability, not production deployment.
Configured timeout is not proof of cancellation under load. Authentication, SERVICE egress,
concurrent writers, memory limits, disaster recovery and production TLS require separate
acceptance tests. The verified read-only mode is access restriction, not a full security certification.

\section{Five-star Linked Open Data assessment}
\label{sec:five-star}
The five-star scheme evaluates how data is published and connected, not the overall
quality of a program or the correctness of every ontology axiom~\cite{five}. There is no
Fuseki-specific requirement: RDFLib or static RDF publication can support the relevant
mechanisms. Table~\ref{tab:five-stars} assesses the project's data publication rather than
assigning an unsupported software-quality score.

\begin{table}[htbp]\centering\small
\caption{Evidence against the cumulative five-star publication scheme.}\label{tab:five-stars}
\begin{tabularx}{\textwidth}{p{1cm}p{3.8cm}X}\toprule
Level & Criterion, summarized & Project evidence and qualification \\\midrule
1 & Publicly accessible, with reuse terms & Public sample accessible; CC BY-SA 4.0 covers project-controlled compilation and material. Upstream rights remain source-specific; the Ministry source has no stated machine-readable open-data licence. \\
2 & Structured, machine-readable & JSON records and parsed RDF graphs, not only screen images or PDFs. \\
3 & Non-proprietary format & JSON, Turtle and JSON-LD artifacts. \\
4 & Web identifiers and open Semantic Web standards & Stable-registry HTTP IRIs, RDF descriptions and query interfaces; sampled entity documents parse. Canonical negotiation remains static; version artifacts are now implemented on the review branch. \\
5 & Connections to other datasets & Corrected local link graph contains 2,309 \term{owl:sameAs}, 61 \term{skos:closeMatch} and 1,188 document-topic links. Existence is measured; every identity assertion is not independently verified. \\
\bottomrule\end{tabularx}
\end{table}

\textbf{Verdict.} The implementation and sampled public RDF support the five-star
publication pattern, subject to the declared licensing scope. Commit \term{a977a239} was
observed live on Render on 9 October 2026 with the corrected 87,183-triple release. This is
not an independent certification, proof of every fact, or proof of a production-grade SLA.
Temporal-evidence and operations limits must therefore accompany the claim rather than be
hidden by a five-star badge.

\subsection{Dated HTTP evidence}
The public Render deployment was checked after manual promotion of commit
\term{a977a239} on 9 October 2026. The checks used bounded read-only HTTP requests against
\url{https://vnedu-lod.onrender.com}; this is an acceptance sample, not continuous
monitoring or a random workload.

\begin{table}[htbp]\centering\small
\caption{Observed HTTP behavior of the deployed corrected release.}\label{tab:http-check}
\begin{tabularx}{\textwidth}{p{4.5cm}X}\toprule
Request & Observation \\\midrule
Render \term{/healthz} & HTTP 200; status \term{ok}, backend \term{rdflib (persistent isolated worker)}, commit \term{a977a239}. \\
Three full triple-count queries & All returned 87,183; observed response times were 11,313, 8,742 and 8,421 ms. \\
Institutional-membership query & Four distinct HUS/HCMUS/USSH results had the intended Hanoi or Ho Chi Minh City parent and an inferred inverse \term{hasMember}; HTTP 200 in 3,477 ms. \\
HUST entity, Accept Turtle & HTTP 200, \term{text/turtle}, 13,347 bytes, 2,653 ms. \\
SPARQL Update attempt & HTTP 400; public query surface remained read-only. \\
Unapproved federated \term{SERVICE} target & HTTP 400; the SSRF allow-list rejected the target. \\
Latest local Fuseki/TDB2 E2E & 87,183 triples loaded, queried, restarted and graph-isomorphism checked; writes rejected with HTTP 405. \\
\bottomrule\end{tabularx}
\end{table}

The Render deployment took 5 minutes 27 seconds and generated the static site during startup;
this is operationally inefficient. The free service can cold-start, uses one bounded query
worker, and the measured 8--11 second full counts are not a high-load benchmark or SLO.
The public runtime uses RDFLib, not the separately accepted Fuseki/TDB2 topology. The probes
do not establish sustained capacity, admin security, backup recovery, or external penetration
test results. They do establish sampled public availability, deployed-revision identity,
read-only behavior and query correctness for the stated release.

The media distinction also matters. For example, the shipped HUST logo metadata records
``fair use'' while its campus photograph records CC0. Describing or linking such an image
does not relicense the image under the dataset license. The revised publisher shows a source link, not an embedded image, for unknown, fair-use or unrecognized licenses. A data publication claim must not
be expanded into a blanket claim that every remotely displayed asset is openly licensed.
